\documentclass[conference]{IEEEtran}
\usepackage[spanish]{babel}
\usepackage[utf8]{inputenc}
\usepackage[T1]{fontenc}
\usepackage{booktabs}
\usepackage{url}

\title{Representacion Tabular de Arboles Binarios}
\author{\IEEEauthorblockN{Ronquillo Nunez Braulio}
\IEEEauthorblockA{Analisis y Diseno de Algoritmos}}

\begin{document}
\maketitle

\begin{abstract}
Los arboles son estructuras jerarquicas fundamentales en algoritmos y
estructuras de datos. Un arbol binario restringe cada nodo a tener a lo mas
dos hijos, lo que permite definir recorridos recursivos y representaciones
compactas. Este reporte describe el modelo teorico de arbol, su
representacion mediante tabla de nodos, los recorridos clasicos y la
complejidad del programa desarrollado.
\end{abstract}

\begin{IEEEkeywords}
arbol binario, recorrido, preorder, inorder, postorder, estructura de datos.
\end{IEEEkeywords}

\section{Introduccion}
Un arbol permite representar informacion con relaciones padre-hijo: sistemas
de archivos, expresiones aritmeticas, decisiones, jerarquias y busquedas. A
diferencia de una lista, la estructura no es lineal. NIST define un arbol como
una estructura que se accede desde una raiz, donde cada nodo puede ser hoja o
nodo interno con hijos \cite{nist_tree}. Tambien puede verse como un grafo
conexo, no dirigido y aciclico.

\section{Marco Teorico}
En un arbol binario, cada nodo tiene como maximo dos hijos: izquierdo y
derecho. La distincion entre ambos importa porque el orden de los hijos afecta
los recorridos y, en otros contextos, la busqueda. Una definicion recursiva
simple establece que un arbol es vacio o esta formado por una raiz y dos
subarboles.

Los recorridos principales son:
\begin{itemize}
\item \textit{Preorder}: visita raiz, subarbol izquierdo y subarbol derecho.
\item \textit{Inorder}: visita subarbol izquierdo, raiz y subarbol derecho.
\item \textit{Postorder}: visita subarbol izquierdo, subarbol derecho y raiz.
\end{itemize}
Estos recorridos visitan cada nodo exactamente una vez. La pagina de notas de
MIT OCW para algoritmos incluye arboles binarios como tema de estructuras de
datos y recorridos \cite{mit_algorithms}.

\section{Representacion Tabular}
El programa usa una tabla con cuatro campos:
\begin{table}[htbp]
\caption{Campos de cada nodo}
\centering
\begin{tabular}{ll}
\toprule
Campo & Significado \\
\midrule
\texttt{id} & identificador unico del nodo \\
\texttt{valor} & dato almacenado \\
\texttt{izq} & identificador del hijo izquierdo \\
\texttt{der} & identificador del hijo derecho \\
\bottomrule
\end{tabular}
\end{table}

El valor $-1$ indica ausencia de hijo. Con esta informacion se construye un
diccionario que permite consultar cualquier nodo por identificador. Despues,
desde la raiz indicada, se ejecutan los recorridos.

\section{Validacion}
La representacion debe cumplir condiciones minimas:
\begin{enumerate}
\item todos los identificadores referenciados deben existir,
\item no debe aparecer un ciclo,
\item un nodo no debe quedar conectado como hijo de varios padres,
\item la raiz debe permitir alcanzar los nodos del arbol.
\end{enumerate}
Estas reglas evitan que una tabla arbitraria se interprete como arbol cuando
en realidad describe un grafo invalido.

\section{Complejidad}
Construir la tabla cuesta $O(n)$ para $n$ nodos. Cada recorrido tambien cuesta
$O(n)$ porque visita todos los nodos una sola vez. La memoria adicional se
debe al diccionario de nodos y a la profundidad de recursion. En un arbol
desbalanceado, la profundidad puede ser $O(n)$; en uno balanceado, $O(\log n)$.

\section{Conclusion}
La representacion tabular es util para observar explicitamente los enlaces
entre nodos. Aunque no usa punteros reales, conserva la idea de raiz e hijos.
Los recorridos preorder, inorder y postorder muestran como una definicion
recursiva se transforma en algoritmos concretos.

\begin{thebibliography}{00}
\bibitem{nist_tree} P. E. Black and Algorithms and Theory of Computation Handbook, ``tree,'' Dictionary of Algorithms and Data Structures, NIST, 2017. [Online]. Available: \url{https://www.nist.gov/dads/HTML/tree.html}
\bibitem{mit_algorithms} E. Demaine, J. Ku, and J. Solomon, ``6.006 Introduction to Algorithms: Lecture Notes,'' MIT OpenCourseWare, 2020. [Online]. Available: \url{https://ocw.mit.edu/courses/6-006-introduction-to-algorithms-spring-2020/pages/lecture-notes/}
\bibitem{ods} P. Morin, ``Open Data Structures,'' Athabasca University Press, 2013. [Online]. Available: \url{https://opendatastructures.org/}
\end{thebibliography}

\end{document}
